\subsection{II. Dạng 1: Xác định vận tốc và quãng đường}
Phương pháp: Xác định các đại lượng đã biết ($v_0, a, t$), sử dụng phương trình động học thích hợp để tìm đại lượng chưa biết.

\begin{vd}
Một ô tô chuyển động nhanh dần đều từ trạng thái nghỉ với gia tốc $a = 2\ \mathrm{m/s^2}$. Tính vận tốc sau $5\ \mathrm{s}$.
\loigiai{
Vận tốc ban đầu $v_0 = 0$. Vận tốc sau $5\ \mathrm{s}$ là:
$v = v_0 + at = 0 + 2 \times 5 = 10\ \mathrm{m/s}$.
}
\end{vd}

% Liên kết đến bài tập vận dụng
\subsubsection{Bài tập tự luyện dạng 1}

\begin{ex}
Một đoàn tàu đang chạy với vận tốc $36\ \mathrm{km/h}$ thì hãm phanh chuyển động chậm dần đều và dừng lại sau $10\ \mathrm{s}$. Gia tốc của tàu là:
\choice
{$-1\ \mathrm{m/s^2}$}
{$1\ \mathrm{m/s^2}$}
{$-3{,}6\ \mathrm{m/s^2}$}
{$3{,}6\ \mathrm{m/s^2}$}
\loigiai{
Đổi $v_0 = 36\ \mathrm{km/h} = 10\ \mathrm{m/s}$. Khi dừng lại, $v = 0$.
Gia tốc: $a = \frac{v - v_0}{t} = \frac{0 - 10}{10} = -1\ \mathrm{m/s^2}$.
}
\end{ex}

\begin{ex}
Một vật bắt đầu chuyển động thẳng nhanh dần đều với gia tốc $a = 0{,}5\ \mathrm{m/s^2}$. Tính quãng đường vật đi được trong $8\ \mathrm{s}$.
\shortans{16 m}
\loigiai{
Quãng đường đi được: $s = \frac{1}{2}at^2 = \frac{1}{2} \cdot 0{,}5 \cdot 8^2 = 16\ \mathrm{m}$.
}
\end{ex}

